\documentclass[11pt]{article}
\usepackage[T1]{fontenc}
\usepackage[utf8]{inputenc}
\usepackage{lmodern}
\usepackage{amsmath,amssymb,amsthm,mathtools}
\usepackage{array,booktabs}
\usepackage{microtype}
\usepackage[margin=1in]{geometry}
\usepackage[hidelinks]{hyperref}
\usepackage{xcolor}

\newtheorem{theorem}{Theorem}[section]
\newtheorem{proposition}[theorem]{Proposition}
\newtheorem{corollary}[theorem]{Corollary}
\newtheorem{remark}[theorem]{Remark}
\newcommand{\C}{\mathbb C}
\newcommand{\R}{\mathbb R}
\newcommand{\N}{\mathbb N}
\newcommand{\SigmaL}{\Sigma_L}
\newcommand{\ModuleHash}{\texttt{199229725d2d359e19a08792317d7f7e552cf027990505a4a4ffb2542295437a}}
\newcommand{\OleanHash}{\texttt{2741f4c24178f2ea632d914f309acc71c714b831a5c5b282b9e8bd8a94a34478}}
\newcolumntype{P}[1]{>{\raggedright\arraybackslash}p{#1}}

\title{Exact Intertwining and Uniform Clustering in the\\
Reconstructed Site Form for Coupled $\mathbb Z_2$ Slices}
\author{Llu\'is Eriksson}
\date{4 August 2026}

\begin{document}
\maketitle

\begin{abstract}
For a finite reflection-positive coupled $\mathbb Z_2$ slice, we identify the
connected two-point function written directly with the reconstructed site form
and the forced transfer operator with the standard connected correlator of a
symmetric Euclidean matrix operator.  The identification is exact at every
time separation.  It follows that the already established Dobrushin
fluctuation bound gives
\[
  |C^{\mathrm{rec}}_L(u;n)|\leq \|u\|_2^2 e^{-mn}
\]
with one $m>0$ chosen before the spatial extent $L$.  The proof preserves both
the exponent and the prefactor.  Nine interface theorems are elaborated by
Lean~4 against Mathlib, including the finite-iterate intertwining, the exact
correlator identity, and the reconstructed-site clustering endpoint.  The new
contribution is this operator/correlator composition.  The small-coupling
Dobrushin estimate is inherited unchanged; no new analytic estimate or
continuum reconstruction is claimed.
\end{abstract}

\section{Result and claim boundary}

Reflection positivity supplies a positive form on boundary data and makes a
time-translation operator meaningful in the reconstructed space
\cite{OS1,OS2,OsterwalderSeiler}.  A spectral estimate written for an auxiliary
Euclidean matrix, however, is not yet a bound on the correlator expressed in
that physical form.  One must identify the two coordinate descriptions,
transport all finite powers, and check that the vacuum subtraction agrees.
This paper carries out exactly those steps for the coupled $\mathbb Z_2$ slice.

The endpoint is stronger than an operator-only interface: the reconstructed
connected correlator itself obeys a volume-uniform exponential estimate.  It
is also deliberately narrower than a general Osterwalder--Schrader (OS)
reconstruction theorem.  We do not prove a new contraction estimate, a
thermodynamic or continuum limit, or a Wightman reconstruction.  The analytic
input is a previously checked finite-state Dobrushin theorem.  The present
work proves that its gap controls the correlator defined by the previously
checked reflection-positive site form, without a change of constants.

\section{Three descriptions of the same finite slice}

For $L\in\N$, let
\[
  \SigmaL=(\operatorname{Fin} L\to\operatorname{Fin}2)
\]
be the finite set of spin configurations.  Fix a strictly positive weight
$w:\SigmaL\to\R_{>0}$ and the real symmetric spatial kernel
$K_\beta(\sigma,\tau)$.  The reconstructed site form on complex boundary
vectors is
\begin{equation}\label{eq:site-form-v2}
  \langle f,g\rangle_{\mathrm{site},w}
  =\sum_{\sigma\in\SigmaL}
    \overline{f(\sigma)}g(\sigma)w(\sigma)^{-1}.
\end{equation}
The transfer operator forced by the site--bond identity is
\begin{equation}\label{eq:forced-v2}
  (T_{w,\beta}g)(\sigma)
  =w(\sigma)\sum_{\tau\in\SigmaL}K_\beta(\sigma,\tau)g(\tau).
\end{equation}
The symmetric weighted matrix used in the spectral lane is
\begin{equation}\label{eq:sym-v2}
  S_{w,\beta}(\sigma,\tau)
  =\sqrt{w(\sigma)}K_\beta(\sigma,\tau)\sqrt{w(\tau)}.
\end{equation}
These objects were already present in the formal development.  The bridge is
the square-root dressing
\begin{equation}\label{eq:q-v2}
  (Q_wu)(\sigma)=\sqrt{w(\sigma)}u(\sigma),
  \qquad u:\SigmaL\to\R,
\end{equation}
with the result regarded as a complex boundary vector.

\begin{proposition}[Isometry and one-step intertwining]\label{prop:intertwine}
If $w(\sigma)>0$ for every $\sigma$, then for real $u,v$,
\begin{align}
 \langle Q_wu,Q_wv\rangle_{\mathrm{site},w}
   &=\sum_\sigma u(\sigma)v(\sigma),\label{eq:isom-v2}\\
 T_{w,\beta}(Q_wu)&=Q_w(S_{w,\beta}u).\label{eq:one-step-v2}
\end{align}
Consequently, for every $n\in\N$,
\begin{equation}\label{eq:iterate-v2}
  T_{w,\beta}^{,n}(Q_wu)=Q_w(S_{w,\beta}^{,n}u).
\end{equation}
\end{proposition}

The equality \eqref{eq:isom-v2} cancels
$\sqrt{w(\sigma)}^2/w(\sigma)$ pointwise.  Equation
\eqref{eq:one-step-v2} is the exact conjugation of the forced transfer
operator, not a comparison in norm.  Equation \eqref{eq:iterate-v2} follows by
induction, so no diagonalisation or unrecorded spectral assumption enters.

Let $\lambda\in\R$ and set
\begin{equation}\label{eq:tilt-v2}
 P_{w,\beta,\lambda}(\sigma,\tau)
 =\lambda^{-1}S_{w,\beta}(\sigma,\tau).
\end{equation}
Scalar multiplication on the complex boundary space and matrix action on the
real Euclidean space obey
\begin{equation}\label{eq:tilt-iterate-v2}
 \bigl(\lambda^{-1}T_{w,\beta}\bigr)^n(Q_wu)
 =Q_w\bigl(P_{w,\beta,\lambda}^{,n}u\bigr)
 \qquad(n\in\N).
\end{equation}
The formal statement uses function iteration for the bare matrix action; a
separate lemma identifies it exactly with powers of the continuous linear
operator $\operatorname{opOf}(P)$ on finite Euclidean space.

\section{The reconstructed connected correlator}

Let $\Omega:\SigmaL\to\R_{>0}$ be a Perron vector and write
\begin{equation}\label{eq:vac-v2}
 \widehat\Omega(\sigma)
 =\frac{\Omega(\sigma)}
 {\sqrt{\sum_\tau\Omega(\tau)^2}}.
\end{equation}
For a real boundary observable $u$ and a time separation $n$, define directly
in the reconstructed site coordinates
\begin{align}
 C^{\mathrm{rec}}_{w,\beta,\lambda,\Omega}(u;n)
  :={}&\operatorname{Re}\left\langle Q_wu,
      (\lambda^{-1}T_{w,\beta})^n Q_wu
       \right\rangle_{\mathrm{site},w}\notag\\
 &-\operatorname{Re}\left\langle Q_w\widehat\Omega,
       Q_wu\right\rangle_{\mathrm{site},w}^{,2}.
 \label{eq:rec-corr}
\end{align}
This is the connected, or truncated, two-point function: the second term is
the squared vacuum expectation.  Both terms use the same reconstructed form
and the same dressing map.

The spectral lane uses the real Euclidean correlator
\begin{equation}\label{eq:conn-standard}
 C_P(u;n)=\langle u,(\operatorname{opOf}P)^n u\rangle_2
   -\langle\widehat\Omega,u\rangle_2^2.
\end{equation}
The central new identity closes the gap between \eqref{eq:rec-corr} and
\eqref{eq:conn-standard}.

\begin{theorem}[Exact correlator identification]\label{thm:corr-id}
For every positive $w$, every real $\beta,\lambda$, every real $\Omega,u$, and
every $n\in\N$,
\begin{equation}\label{eq:corr-id}
 C^{\mathrm{rec}}_{w,\beta,\lambda,\Omega}(u;n)
 =C_{P_{w,\beta,\lambda}}(u;n).
\end{equation}
\end{theorem}

The proof applies \eqref{eq:tilt-iterate-v2} to the time-dependent term,
\eqref{eq:isom-v2} to both pairings, and the matrix-power lemma to change from
iteration to $\operatorname{opOf}(P)^n$.  It is an equality for every finite
$n$; there is no limiting argument and no loss of exponent or prefactor.

\section{Gap implies reconstructed-site clustering}

Let
\[
 \Pi_\Omega v=\langle\widehat\Omega,v\rangle_2\widehat\Omega,
 \qquad P_\perp=\operatorname{opOf}(P)-\Pi_\Omega.
\]
If $P$ is symmetric, $P\Omega=\Omega$, and $\Omega$ is positive, the formal
library packages $\operatorname{opOf}(P)$ and $\widehat\Omega$ as a vacuum
transfer pair.  The existing gap-to-clustering theorem gives
\begin{equation}\label{eq:gap-cluster}
 |C_P(u;n)|\leq \|u\|_2^2 r^n
 \quad\text{whenever}\quad \|P_\perp\|\leq r.
\end{equation}
Combining this with Theorem~\ref{thm:corr-id} yields the physical-form
endpoint.

\begin{corollary}[Reconstructed-site decay]\label{cor:rec-decay}
Under the hypotheses above, if $\|P_\perp\|\leq r$, then
\begin{equation}\label{eq:rec-decay}
 |C^{\mathrm{rec}}_{w,\beta,\lambda,\Omega}(u;n)|
 \leq \|u\|_2^2 r^n
\end{equation}
for every real boundary vector $u$ and every $n\in\N$.
\end{corollary}

The result preserves the exact constant in \eqref{eq:gap-cluster}.  In
particular, it does not replace $r$ by a larger comparison rate or introduce a
volume-dependent equivalence-of-norms factor.

\section{The volume-uniform endpoint}

For the coupled-slice weight $w=\operatorname{sliceW}(\gamma,L)$, the inherited
Dobrushin theorem supplies a positive exponent before the spatial size is
chosen.  The formal convention uses configurations on $\Sigma_{L+1}$; we keep
that binder rather than silently renumbering it.

\begin{theorem}[Uniform reconstructed-site clustering]\label{thm:uniform-v2}
Let $0<\alpha<1$ and suppose
\begin{equation}\label{eq:window-v2}
 2\tanh|\beta|+2\tanh|\gamma|\leq\alpha.
\end{equation}
Then there exists $m>0$ such that, for every $L\in\N$, there are a positive
Perron value $\lambda_L$ and a positive vector
$\Omega_L:\Sigma_{L+1}\to\R_{>0}$ satisfying
\begin{align}
 P_L\Omega_L&=\Omega_L,\label{eq:fixed-v2}\\
 \left\|\operatorname{opOf}(P_L)-\Pi_{\Omega_L}\right\|
   &\leq e^{-m},\label{eq:gap-v2}
\end{align}
where
$P_L=P_{\operatorname{sliceW}(\gamma,L),\beta,\lambda_L}$.  Moreover, for all
real $u:\Sigma_{L+1}\to\R$ and all $n\in\N$,
\begin{equation}\label{eq:uniform-corr-v2}
 \left|C^{\mathrm{rec}}_{
   \operatorname{sliceW}(\gamma,L),\beta,\lambda_L,\Omega_L}(u;n)\right|
 \leq \|u\|_2^2 e^{-mn},
\end{equation}
and the exact intertwining \eqref{eq:tilt-iterate-v2} holds.
\end{theorem}

The quantifier order
\[
 \exists m>0\;\forall L\;\exists(\lambda_L,\Omega_L)
\]
is part of the theorem type.  This, rather than the mere existence of a
finite-volume Perron gap for each $L$, is the asserted uniformity.  The new
clause \eqref{eq:uniform-corr-v2} makes the endpoint a statement about the
reconstructed connected correlator, not only an auxiliary operator norm.

\section{Dependency ledger: inherited and proved here}

The proof has a short exact layer on top of a longer analytic dependency.
Keeping them separate prevents a compositional result from being advertised
as a new estimate.

\begin{center}
\small
\begin{tabular}{@{}P{0.27\linewidth}P{0.30\linewidth}P{0.34\linewidth}@{}}
\toprule
Layer & Status in this paper & Mathematical role \\
\midrule
Reflection-positive site form and forced $T_{w,\beta}$ & inherited & identifies the reconstructed finite boundary space and its time translation \\
Symmetric weighted kernel and Perron tilt & inherited & supplies the real Euclidean operator used by the spectral lane \\
Dobrushin window \eqref{eq:window-v2} and bound \eqref{eq:gap-v2} & inherited unchanged & chooses $m$ before $L$ \\
Square-root isometry and all-iterate intertwining & proved in the interface module & identifies reconstructed and Euclidean coordinates exactly \\
Matrix powers versus bare action iteration & proved in the interface module & aligns the two implementations of finite time evolution \\
Identity \eqref{eq:corr-id} and decay \eqref{eq:rec-decay} & proved in the interface module & transports the gap to the reconstructed connected correlator without loss \\
Uniform endpoint \eqref{eq:uniform-corr-v2} & proved by composition & exposes the physical-form conclusion with the correct binder order \\
\bottomrule
\end{tabular}
\end{center}

This ledger also describes the paper's autonomy.  The result depends on
machine-checked repository modules, but its new claim is not merely that those
files coexist: the exact correlator equality and its uniform corollary are
new theorem declarations whose types mention both lanes.

\section{Formal artifact and reproducibility}

The authoritative file is
\texttt{YangMills/OS/OSReconstructionUniform.lean}.  It contains two
definitions and nine theorem declarations:

\begin{center}
\small
\begin{tabular}{@{}P{0.40\linewidth}P{0.50\linewidth}@{}}
\toprule
Lean declaration & Role \\
\midrule
\texttt{siteForm\_qEmbed} & reconstructed-form isometry \\
\texttt{transferOp\_qEmbed} & one-step forced-operator intertwining \\
\texttt{transferOp\_iterate\_qEmbed} & unnormalised finite iterates \\
\texttt{transferOp\_qEmbed\_tilt} & one-step Perron-normalised intertwining \\
\texttt{transferOp\_qEmbed\_tilt\_iterate} & normalised finite iterates \\
\texttt{opOf\_pow\_toLp\_act} & operator powers equal matrix-action iteration \\
\texttt{reconstructedConnCorr\_eq\_connCorr} & exact identity \eqref{eq:corr-id} \\
\texttt{reconstructedConnCorr\_decay} & reconstructed-site decay \eqref{eq:rec-decay} \\
\texttt{os\_reconstruction\_uniform\_gap} & uniform theorem, including \eqref{eq:uniform-corr-v2} \\
\bottomrule
\end{tabular}
\end{center}

Validation used a Colab Pro+ CPU/high-RAM runtime with the GPU disabled; no
Lean or Lake command ran on the local Windows desktop.  The target build
completed all 8191 jobs with exit code zero.  The exact LF source had 8682
bytes and SHA-256
\[
 \ModuleHash,
\]
while the materialised \texttt{.olean} had SHA-256
\[
 \OleanHash.
\]
A separate verification manifest records timestamps, toolchain pins, oracle
counters, byte counts, and attacks attempted.  Compilation and counters are
evidence for an independent audit, not a self-issued terminal verdict.

\section{Relation to primary literature}

The OS papers formulate Euclidean reconstruction and reflection positivity
\cite{OS1,OS2}; Osterwalder and Seiler develop the lattice-gauge setting
\cite{OsterwalderSeiler}; and Fr\"ohlich, Israel, Lieb, and Simon use reflection
positivity in lattice systems \cite{FILS}.  Dobrushin's conditional-probability
criterion is the primary source for the high-temperature uniqueness mechanism
\cite{Dobrushin1968}.  These are conceptual and analytic antecedents, not
claims of novelty.

The closest public papers in the author's own sequence separate the
reflection-positive site form, the transfer bridge, the finite-volume Perron
vacuum, and fluctuation-sector bounds.  The present paper's distinct endpoint
is the exact equality \eqref{eq:corr-id} followed by the reconstructed-site
bound \eqref{eq:uniform-corr-v2}.  We make no global priority claim: unindexed
formal developments may exist, and bibliographic search is not a mathematical
proof of precedence.

\section{Limitations}

The state space is finite and specific to coupled $\mathbb Z_2$ slices.  The
paper does not prove a general OS theorem, a continuum field theory, or a
Yang--Mills mass gap.  It does not improve the Dobrushin region
\eqref{eq:window-v2} or compute a sharper $m$.  The observable is a real
boundary vector; extensions to larger observable algebras are not addressed.
Finally, \eqref{eq:uniform-corr-v2} is an exponential time-separation estimate
inside the reconstructed finite-volume site space.  Calling it more than that
would exceed the formal theorem.

\section{Conclusion}

The square-root dressing does more than match two transfer matrices: it
matches their iterates, vacuum expectation, and connected correlator.  That
identity turns the inherited uniform fluctuation bound into an explicit
reconstructed-site clustering theorem with the same rate and prefactor.  The
result is a compact but nontrivial endpoint whose dependency, quantifier order,
and limitations are all visible in the machine-checked statement.

\begin{thebibliography}{9}

\bibitem{OS1}
K.~Osterwalder and R.~Schrader,
\newblock Axioms for Euclidean Green's functions,
\newblock \emph{Communications in Mathematical Physics} \textbf{31} (1973), 83--112.
\newblock \href{https://doi.org/10.1007/BF01645738}{doi:10.1007/BF01645738}.

\bibitem{OS2}
K.~Osterwalder and R.~Schrader,
\newblock Axioms for Euclidean Green's functions II,
\newblock \emph{Communications in Mathematical Physics} \textbf{42} (1975), 281--305.
\newblock \href{https://doi.org/10.1007/BF01608978}{doi:10.1007/BF01608978}.

\bibitem{OsterwalderSeiler}
K.~Osterwalder and E.~Seiler,
\newblock Gauge field theories on a lattice,
\newblock \emph{Annals of Physics} \textbf{110} (1978), 440--471.
\newblock \href{https://doi.org/10.1016/0003-4916(78)90039-8}{doi:10.1016/0003-4916(78)90039-8}.

\bibitem{FILS}
J.~Fr\"ohlich, R.~Israel, E.~H.~Lieb, and B.~Simon,
\newblock Phase transitions and reflection positivity. I. General theory and long range lattice models,
\newblock \emph{Communications in Mathematical Physics} \textbf{62} (1978), 1--34.
\newblock \href{https://doi.org/10.1007/BF01940327}{doi:10.1007/BF01940327}.

\bibitem{Dobrushin1968}
R.~L.~Dobrushin,
\newblock The description of a random field by means of conditional probabilities and conditions of its regularity,
\newblock \emph{Theory of Probability and its Applications} \textbf{13} (1968), 197--224.
\newblock \href{https://doi.org/10.1137/1113026}{doi:10.1137/1113026}.

\end{thebibliography}

\end{document}
